\subsection*{الف}
فاصلهٔ اقلیدسی ($q=2$) تنها فرم دورانی و انتقالی‌ناوردا است.

ناوردایی انتقالی برای همهٔ $q$ برقرار است، زیرا فاصله تنها به تفاضل $x_k - y_k$ وابسته است:
\begin{flalign*}
    L_q(x+t,\, y+t) &= \left(\sum_{k=1}^{d} |(x_k+t_k) - (y_k+t_k)|^q\right)^{1/q} = L_q(x,y) &&
\end{flalign*}

اما ناوردایی دورانی تنها برای $q=2$ برقرار است. برای ماتریس دوران $R$ ($R^\top R = I$):
\begin{flalign*}
    L_2(Rx, Ry)^2 &= \|R(x-y)\|^2 = (x-y)^\top R^\top R (x-y) \\
    &= (x-y)^\top (x-y) = \|x-y\|^2 = L_2(x,y)^2 &&
\end{flalign*}
علت بنیادین این است که نُرم اقلیدسی از یک \lr{inner product} القا می‌شود و دوران‌ها همان تبدیل‌های متعامد حافظِ این \lr{inner product} اند؛ سایر مقادیر $q$ چنین ساختاری ندارند. برای مثال در $q=1$ با دوران $45^\circ$:
\begin{flalign*}
    x-y = (1,0) &\implies L_1 = 1 \\
    R_{45^\circ}(x-y) = \left(\tfrac{1}{\sqrt 2}, \tfrac{1}{\sqrt 2}\right) &\implies L_1 = \sqrt 2 \neq 1 &&
\end{flalign*}

\subsection*{ب}
فاصلهٔ \lr{Manhattan} ($q=1$) هنگامی مناسب است که حرکت تنها در راستای محورها مجاز باشد و حرکت قطری معنا نداشته باشد؛ مانند مسیریابی در شبکهٔ خیابان‌های شطرنجی یا حرکت روی \lr{grid}. همچنین در فضاهای با ابعاد بالا به دلیل حساسیت کمتر به \lr{outlier}ها (وزن‌دهی خطی به تفاوت‌ها به‌جای درجهٔ دوم) بر اقلیدسی ترجیح داده می‌شود، و در مواردی که ویژگی‌ها مستقل و ناهمگن‌اند تفسیرپذیری بهتری دارد.